% ============================================================
% CHƯƠNG 4: KIẾN TRÚC YOLO26
% ============================================================
\chapter{Kiến trúc mô hình YOLO26}
\label{ch:architecture}

Chương này trình bày chi tiết kiến trúc mô hình \yolo, bao gồm ba phần chính: Backbone (trích xuất đặc trưng), Neck (tổng hợp đặc trưng đa tỷ lệ), và Detection Head (dự đoán bounding box và lớp đối tượng). Nội dung được tổng hợp từ paper gốc~\cite{yolo26}, mã nguồn Ultralytics~\cite{ultralytics}, và quá trình nghiên cứu xây dựng lại mô hình từ đầu (from scratch) của nhóm.

\begin{figure}[H]
    \centering
    \includegraphics[width=\textwidth]{YOLO26.drawio.png}
    \caption{Sơ đồ khối tổng quát kiến trúc YOLO26. Đầu vào ảnh $3 \times 640 \times 640$ được xử lý qua Backbone (trích xuất đặc trưng), Neck (tổng hợp FPN+PAN), và Detection Head (Dual-Head O2M/O2O).}
    \label{fig:yolo26_overview}
\end{figure}

% ===================================================================
\section{Backbone}
\label{sec:arch_backbone}

Backbone của \yolo gồm 11 lớp (b0--b10) có nhiệm vụ trích xuất đặc trưng ảnh ở nhiều mức phân giải khác nhau.

\subsection{Khối Conv cơ bản}

Trong \yolo, mỗi khối \texttt{Conv} bao gồm ba phần: lớp tích chập (\texttt{Conv2d}), chuẩn hóa batch (\texttt{BatchNorm2d}), và hàm kích hoạt \texttt{SiLU}:

\begin{lstlisting}[caption={Khối Conv cơ bản trong YOLO26}, label=lst:conv]
class Conv(nn.Module):
    def __init__(self, c1, c2, k=1, s=1, p=None, g=1):
        self.conv = nn.Conv2d(c1, c2, k, s, autopad(k, p), groups=g, bias=False)
        self.bn = nn.BatchNorm2d(c2)
        self.act = nn.SiLU(inplace=True)
    
    def forward(self, x):
        return self.act(self.bn(self.conv(x)))
\end{lstlisting}

Các lớp Conv đầu tiên (b0, b1, b3, b5, b7) sử dụng kernel $3 \times 3$ với stride $s=2$ để giảm kích thước không gian xuống một nửa ở mỗi bước.

\textbf{Conv 1x1} trộn channel tại cùng vị trí spatial. \textbf{Conv 3x3} trộn thông tin không gian (cạnh, biên, texture).

\subsection{Khối DWConv (Depthwise Convolution)}
\label{subsec:dwconv}

DWConv là tích chập tách biệt theo chiều sâu (Depthwise Convolution), trong đó số lượng nhóm tính toán bằng số kênh vào và số kênh ra ($\text{groups} = \gcd(c_1, c_2)$):
\begin{equation}
    \text{Params}_{\text{DW}} = c_1 \times k^2 \quad \ll \quad \text{Params}_{\text{std}} = c_1 \times c_2 \times k^2
\end{equation}
DWConv được sử dụng rộng rãi trong classification head (cv3) của YOLO26 nhằm giảm chi phí tính toán và tăng tốc độ suy luận.

\subsection{Khối Bottleneck}
\label{subsec:bottleneck}

Bottleneck là block residual cơ bản, tạo đặc trưng $F(x)$ bằng 2 lớp Conv rồi tùy điều kiện cộng residual với đầu vào $x$:

\begin{figure}[H]
\centering
\begin{tikzpicture}[node distance=1cm and 1.2cm, auto]
    % Nodes
    \node (input) [splitnode] {};
    \node (lbl_in) [left=0.15cm of input] {Input $x$};
    \node (cv1) [block, right=of input] {Conv 1x1\\(cv1)};
    \node (cv2) [block, right=of cv1] {Conv 3x3 / DWConv\\(cv2)};
    \node (sum) [sum, right=of cv2];
    \node (output) [right=0.8cm of sum] {Output};

    % Connections
    \path [line] (input) -- (cv1);
    \path [line] (cv1) -- (cv2);
    \path [line] (cv2) -- (sum);
    \path [line] (sum) -- (output);
    
    % Skip connection
    \path [line] (input) -- ++(0,1.0) -| (sum.north);
    
    % Labels
    \node [above=0.6cm of cv1, font=\small\itshape] {Skip Connection (Shortcut = True)};
\end{tikzpicture}
\caption{Sơ đồ khối Bottleneck với skip connection.}
\label{fig:bottleneck}
\end{figure}

\begin{equation}
    F(x) = \text{cv2}(\text{cv1}(x))
    \label{eq:bottleneck_Fx}
\end{equation}
\begin{equation}
    \text{Output} = \begin{cases}
        x + F(x) & \text{nếu } \texttt{shortcut} = \texttt{True} \text{ và } c_1 = c_2 \\
        F(x) & \text{ngược lại}
    \end{cases}
    \label{eq:bottleneck_output}
\end{equation}

\noindent trong đó:
\begin{itemize}
    \item \texttt{cv1}: Conv$(c_1, c\_, k_1, s=1)$ với $c\_ = \lfloor c_2 \times e \rfloor$ (kênh ẩn, $e$ là expansion ratio).
    \item \texttt{cv2}: Conv$(c\_, c_2, k_2, s=1)$ với $\texttt{groups}=g$.
    \item Nhánh skip chỉ kích hoạt khi $c_1 = c_2$ và \texttt{shortcut=True}.
\end{itemize}

\textbf{Tác dụng}: Giảm chi phí tính toán (nhờ expansion ratio $e < 1$), cho phép xây dựng mạng sâu hơn mà vẫn nhẹ, và giữ thông tin gốc nhờ nhánh residual giúp tránh vanishing gradient.

Ví dụ: input $[B, 128, 80, 80]$, $e=0.5$ $\to$ $c\_=64$ $\to$ output $[B, 128, 80, 80]$ (chiều dài và chiều rộng giữ nguyên).

\subsection{Khối C2f (Cross Stage Partial with 2 Convolutions)}
\label{subsec:c2f}

Khối C2f~\cite{cspnet} kết hợp tư tưởng của CSPNet nhằm gia tăng số lượng đường truyền gradient, giúp mạng học được nhiều đặc trưng đa dạng hơn:

\begin{equation}
    \text{C2f}(x) = \text{cv2}\Big(\text{Concat}\big[y_0, y_1, m_1(y_1), m_2(m_1(y_1)), \ldots\big]\Big)
\end{equation}

Đầu vào sau khi qua \texttt{cv1} được chia đôi thành 2 nhánh: $[y_0, y_1] = \text{chunk}(\text{cv1}(x), 2)$. Nhánh nông $y_0$ đi thẳng không xử lý sâu thêm, trong khi nhánh $y_1$ được truyền qua một chuỗi các khối Bottleneck $m_i$.

\textbf{Sơ đồ luồng dữ liệu (với n=2)}:
\begin{verbatim}
Input -> cv1 -> split -> [y0, y1]
                              |
                         Bottleneck(y1) -> y2
                              |
                         Bottleneck(y2) -> y3
                              |
              cat(y0, y1, y2, y3) -> cv2 -> Output
\end{verbatim}

\textbf{Ý nghĩa}: Cơ chế split-and-concat này vừa giữ lại các đặc trưng nông thô ($y_0$) vừa trích xuất các đặc trưng ngữ cảnh sâu ($y_3$), giúp gradient lan truyền hiệu quả hơn trong quá trình backward.

\subsection{Khối C3 và C3k}
\label{subsec:c3_c3k}

\textbf{C3} là khối CSP kiểu cũ có hai nhánh chính: nhánh xử lý sâu qua chuỗi Bottleneck và nhánh đi thẳng qua một khối Conv, sau đó ghép lại bằng phép concat và xử lý qua Conv 1x1 cuối.

\textbf{C3k} kế thừa C3 nhưng sử dụng Bottleneck với kernel $k \times k$ tùy chỉnh (thường $k=3$):
\begin{equation}
    \text{C3k}: \text{Bottleneck}_{\text{C3}} \text{ dùng } k=((1,1),(3,3)) \quad \to \quad \text{C3k dùng } k=(3,3)
\end{equation}

\begin{figure}[H]
\centering
\begin{tikzpicture}[node distance=0.8cm and 1.2cm, auto]
    % Nodes
    \node (input) [splitnode] {};
    \node (lbl_in) [left=0.15cm of input] {Input $x$};
    
    % Upper branch (Main branch with Bottlenecks)
    \node (cv1) [block, above right=0.2cm and 1.0cm of input] {Conv 1x1\\(cv1)};
    \node (bottlenecks) [block, right=of cv1] {Chuỗi $n$ khối\\C3k / Bottleneck};
    
    % Lower branch (Shortcut / Skip branch)
    \node (cv2) [block, below right=0.2cm and 1.0cm of input] {Conv 1x1\\(cv2)};
    
    % Concat node
    \node (concat) [grayblock, right=1.2cm of bottlenecks, yshift=-0.7cm] {Concat};
    
    % Final Conv
    \node (cv3) [block, right=of concat] {Conv 1x1\\(cv3)};
    \node (output) [right=0.8cm of cv3] {Output};

    % Connections
    \path [line] (input) -- ++(0,0.7) |- (cv1.west);
    \path [line] (input) -- ++(0,-0.7) |- (cv2.west);
    \path [line] (cv1) -- (bottlenecks);
    \path [line] (bottlenecks.east) -- ++(0.4,0) |- (concat.west);
    \path [line] (cv2.east) -- ++(0.4,0) |- (concat.west);
    \path [line] (concat) -- (cv3);
    \path [line] (cv3) -- (output);
\end{tikzpicture}
\caption{Sơ đồ khối C3k --- C3 với kernel 3x3 sâu hơn.}
\label{fig:c3k}
\end{figure}

Khối C3k có khả năng bao quát thông tin không gian tốt hơn nhờ mở rộng kernel tích chập ở nhánh sâu, giúp học được các đặc trưng biên và cấu trúc hình học tốt hơn, đổi lại chi phí tính toán cao hơn Bottleneck thông thường.

\subsection{Khối C3k2}
\label{subsec:c3k2}

C3k2 là khối trộn đặc trưng chính và quan trọng nhất trong \yolo, kế thừa từ C2f nhưng cho phép tùy biến loại block bên trong:

\begin{equation}
    y_0, y_1 = \text{Split}(\text{cv1}(x))
    \label{eq:c3k2_split}
\end{equation}
\begin{equation}
    z_i = \text{Block}_i(z_{i-1}), \quad z_0 = y_1, \quad i = 1, \ldots, n
    \label{eq:c3k2_blocks}
\end{equation}
\begin{equation}
    \text{Output} = \text{cv2}\big(\text{cat}([y_0, y_1, z_1, \ldots, z_n])\big)
    \label{eq:c3k2_output}
\end{equation}

\begin{figure}[H]
\centering
\begin{tikzpicture}[
    node distance=0cm,
    block/.append style={minimum width=1.6cm, minimum height=0.8cm},
    auto
]
    % --- Node definitions ---
    \node (input)  []                            {Input $x$};
    \node (cv1)    [block, right=1.0cm of input] {Conv 1x1\\(cv1)};
    % Place Concat with explicit gap so shortcut wire has room above
    \node (concat) [grayblock, right=5.5cm of cv1] {Concat};
    \node (cv2)    [block,  right=0.8cm of concat] {Conv 1x1\\(cv2)};
    \node (output) [right=0.8cm of cv2]            {Output};

    % Split point on the wire after cv1
    \coordinate (S) at ($(cv1.east)+(0.55,0)$);

    % Block row, centred below the cv1-to-concat gap
    \coordinate (midgap) at ($(cv1.east)!0.5!(concat.west)$);
    % Lower the block row to y = -1.5cm so we have a clear vertical entry into Concat
    \node (block1) [block, below=1.5cm of midgap, xshift=-1.0cm] {Block 1};
    \node (block2) [block, right=0.7cm of block1]               {Block 2};

    % --- Forward main wire ---
    \path [line] (input) -- (cv1);
    \draw [-]    (cv1.east) -- (S);   % no arrowhead at split

    % --- Shortcut y0: go UP by 1.0cm, then RIGHT to concat.north ---
    % Path-based node placement to ensure label is perfectly centered and above the line
    \coordinate (Sup) at ($(S)+(0,1.0)$);
    \path [line] (S) -- (Sup) -- node[above, midway, font=\small] {Nhánh $y_0$ (đi thẳng)} (Sup -| concat.north) -- (concat.north);

    % --- Branch y1: go DOWN then RIGHT to block1.west ---
    % Single bend: goes down to block1's center line, then horizontally to block1.west
    \path [line] (S) |- (block1.west);
    \node [left=5pt, font=\small] at (S |- block1.west) {$y_1$};

    % --- Block 1 to Block 2 ---
    \path [line] (block1) -- (block2);

    % --- Outputs from blocks to concat.south ---
    % Block 1 output (from mid12) goes up to y = -0.9, then right to concat.south-west, then up
    \coordinate (mid12)  at ($(block1.east)!0.5!(block2.west)$);
    \path [line] (mid12) -- ($(mid12)+(0,0.6)$) -| ($(concat.south)+(-0.35,0)$);

    % Block 2 output (from block2.east) goes horizontally right to concat.south-east's x, then straight up
    \path [line] (block2.east) -| ($(concat.south)+(+0.35,0)$);

    % --- Concat onwards ---
    \path [line] (concat) -- (cv2);
    \path [line] (cv2)    -- (output);
\end{tikzpicture}
\caption{Sơ đồ khối C3k2 --- forward theo C2f, block bên trong tùy cấu hình.}
\label{fig:c3k2}
\end{figure}

Tùy vào 2 tham số boolean \texttt{attn} và \texttt{c3k}, các Block bên trong sẽ được khởi tạo khác nhau:
\begin{enumerate}
    \item \texttt{attn=True}: Block = Bottleneck + PSABlock (Position-Sensitive Attention).
    \item \texttt{attn=False, c3k=True}: Block = C3k (kernel $k \times k$ sâu).
    \item \texttt{attn=False, c3k=False}: Block = Bottleneck tiêu chuẩn.
\end{enumerate}

\begin{table}[H]
\centering
\small
\caption{Sử dụng C3k2 tại các vị trí trong YOLO26.}
\label{tab:c3k2_usage}
\begin{tabular}{lcccl}
\toprule
\textbf{Layer} & \textbf{c3k} & \textbf{attn} & \textbf{Module bên trong} & \textbf{Vị trí} \\
\midrule
b2, b4 & False & False & Bottleneck đơn giản & Backbone đầu \\
b6, b8 & True & False & C3k (3x3 sâu) & Backbone P4, P5 \\
h13--h19 & True & False & C3k & Neck \\
h22 & True & True & Bottleneck + PSABlock & Neck P5 output \\
\bottomrule
\end{tabular}
\end{table}

\textbf{Ví dụ sự thay đổi kích thước tensor}: C3k2($c_1=1024, c_2=256, n=2, \text{c3k}=\text{True}, e=0.5$):
\begin{verbatim}
[B, 1024, 80, 80] -> cv1 -> [B, 256, 80, 80]
    -> split: y0 [B, 128, 80, 80], y1 [B, 128, 80, 80]
    -> C3k(y1) -> y2 [B, 128, 80, 80]
    -> C3k(y2) -> y3 [B, 128, 80, 80]
    -> cat(y0,y1,y2,y3) -> [B, 512, 80, 80]
    -> cv2 -> [B, 256, 80, 80]
\end{verbatim}
Kích thước không gian $H, W$ được giữ nguyên hoàn toàn do tất cả các lớp tích chập bên trong đều sử dụng stride=1 và cơ chế tự động đệm (autopad).

\subsection{Khối SPPF (Spatial Pyramid Pooling Fast)}
\label{subsec:sppf}

SPPF mở rộng vùng nhận cảm ngữ cảnh (receptive field) bằng cách áp dụng liên tiếp 3 lớp MaxPool $5 \times 5$ song song với việc lưu lại kết quả trung gian, mà không làm giảm kích thước không gian ảnh đầu vào:

\begin{figure}[H]
\centering
\begin{tikzpicture}[node distance=0cm, auto]
    % --- Top row: processing chain ---
    \node (input) []                             {Input $x$};
    \node (cv1)   [block,       right=0.8cm of input] {Conv 1x1\\(cv1)};
    \node (mp1)   [orangeblock, right=1.5cm of cv1]   {MaxPool 5x5\\(mp1)};
    \node (mp2)   [orangeblock, right=1.5cm of mp1]   {MaxPool 5x5\\(mp2)};
    \node (mp3)   [orangeblock, right=1.5cm of mp2]   {MaxPool 5x5\\(mp3)};

    % --- Bottom row: Concat well below so drop-lines have clearance ---
    \node (concat) [grayblock, below=2.2cm of mp2]     {Concat};
    \node (cv2)    [block,     right=1.0cm of concat]  {Conv 1x1\\(cv2)};
    \node (output) [right=0.8cm of cv2]                {Output};

    % --- Tap coordinates: midpoints of each gap ---
    \coordinate (T0) at ($(cv1.east)!0.5!(mp1.west)$);
    \coordinate (T1) at ($(mp1.east)!0.5!(mp2.west)$);
    \coordinate (T2) at ($(mp2.east)!0.5!(mp3.west)$);
    \coordinate (T3) at ($(mp3.east)+(0.6,0)$);

    % --- Main chain connections ---
    \path [line] (input) -- (cv1);
    \draw [-]    (cv1.east) -- (T0);
    \path [line] (T0)       -- (mp1.west);
    \draw [-]    (mp1.east) -- (T1);
    \path [line] (T1)       -- (mp2.west);
    \draw [-]    (mp2.east) -- (T2);
    \path [line] (T2)       -- (mp3.west);
    \draw [-]    (mp3.east) -- (T3);

    % --- Vertical drops with two separate horizontal levels to prevent overlaps ---
    \coordinate (L_high) at ($(concat.north)+(0,1.2)$);
    \coordinate (L_low)  at ($(concat.north)+(0,0.6)$);

    \coordinate (Concat1) at ($(concat.north)+(-0.6,0)$);
    \coordinate (Concat2) at ($(concat.north)+(-0.2,0)$);
    \coordinate (Concat3) at ($(concat.north)+(+0.2,0)$);
    \coordinate (Concat4) at ($(concat.north)+(+0.6,0)$);

    \path [line] (T0) -- (T0 |- L_high) -| (Concat1);
    \path [line] (T1) -- (T1 |- L_low)  -| (Concat2);
    \path [line] (T2) -- (T2 |- L_low)  -| (Concat3);
    \path [line] (T3) -- (T3 |- L_high) -| (Concat4);

    % --- Concat to output ---
    \path [line] (concat) -- (cv2);
    \path [line] (cv2)    -- (output);
\end{tikzpicture}
\caption{Sơ đồ khối SPPF --- MaxPool liên tiếp tạo multi-scale receptive field.}
\label{fig:sppf}
\end{figure}

\begin{equation}
    y_0 = \text{cv1}(x), \quad
    y_1 = \text{MaxPool}(y_0), \quad
    y_2 = \text{MaxPool}(y_1), \quad
    y_3 = \text{MaxPool}(y_2)
    \label{eq:sppf}
\end{equation}
\begin{equation}
    \text{Output} = \text{cv2}\big(\text{cat}([y_0, y_1, y_2, y_3])\big)
\end{equation}

\begin{figure}[H]
\centering
\begin{tikzpicture}[scale=0.35]
    % Draw grids/boxes
    % Outer box 13x13 (represented as coordinates from -6.5 to 6.5)
    \draw[thick, draw=red!70!black, fill=red!5, fill opacity=0.4] (-6.5,-6.5) rectangle (6.5,6.5);
    % Middle box 9x9 (from -4.5 to 4.5)
    \draw[thick, draw=orange!80!black, fill=orange!5, fill opacity=0.5] (-4.5,-4.5) rectangle (4.5,4.5);
    % Inner box 5x5 (from -2.5 to 2.5)
    \draw[thick, draw=blue!80!black, fill=blue!5, fill opacity=0.6] (-2.5,-2.5) rectangle (2.5,2.5);
    
    % Draw center pixel
    \draw[fill=black] (-0.5,-0.5) rectangle (0.5,0.5);
    \node[white, font=\tiny] at (0,0) {C};

    % Draw grid lines for illustration
    \draw[step=1cm, gray!30, very thin] (-6.5,-6.5) grid (6.5,6.5);

    % Labels
    \node[blue!80!black, font=\small\bfseries] at (0, 1.2) {$5 \times 5$};
    \node[orange!80!black, font=\small\bfseries] at (0, 3.2) {$9 \times 9$};
    \node[red!70!black, font=\small\bfseries] at (0, 5.2) {$13 \times 13$};
    
    % Legends on the right
    \node[anchor=west, font=\small] at (7.5, 3) {\textcolor{blue!80!black}{\textbf{Lần 1 (mp1)}}: Receptive Field $5 \times 5$};
    \node[anchor=west, font=\small] at (7.5, 1) {\textcolor{orange!80!black}{\textbf{Lần 2 (mp2)}}: Receptive Field $9 \times 9$};
    \node[anchor=west, font=\small] at (7.5, -1) {\textcolor{red!70!black}{\textbf{Lần 3 (mp3)}}: Receptive Field $13 \times 13$};
    \node[anchor=west, font=\small] at (7.5, -3) {\textbf{C}: Điểm ảnh trung tâm (Center Pixel)};
\end{tikzpicture}
\caption{Minh họa sự mở rộng receptive field tương đương $5\times5$, $9\times9$, và $13\times13$ qua 3 lần MaxPool 5x5 liên tiếp.}
\label{fig:sppf_maxpool}
\end{figure}

Mỗi lần MaxPool $5 \times 5$ (với stride $s=1$, padding $p=2$) mở rộng vùng ngữ cảnh thêm 4 pixel. Qua 3 lần lặp, thông tin ngữ cảnh toàn cục được tổng hợp tương đương một cửa sổ MaxPool có kích thước $13 \times 13$, hỗ trợ mô hình nhận diện vật thể ở nhiều kích thước khác nhau.

\textbf{So với SPP}: SPP dùng 3 kernel song song ($5\times5$, $9\times9$, $13\times13$); trong khi SPPF dùng cùng một kernel $5\times5$ nối tiếp --- cho receptive field tương đương nhưng \textbf{giảm đáng kể chi phí tính toán}.

\begin{notebox}[Lưu ý: YOLO26 SPPF có act=False trên cv1]
Khác với SPPF gốc (Ultralytics) dùng \texttt{act=True} (SiLU) trên lớp tích chập cv1, YOLO26 sử dụng \texttt{act=False}. Thiết kế này làm giảm dung lượng biểu diễn phi tuyến tính ngay trước bước max-pooling, giúp ổn định thông tin biên của vật thể lớn.
\end{notebox}

\subsection{Cơ chế Attention (Multi-Head Self-Attention)}
\label{subsec:attention}

Trong các cấu hình nâng cao của YOLO26, Multi-Head Self-Attention được chèn vào Neck (h22) và Backbone (b10) để bắt cặp các tương quan không gian toàn cục:
\begin{figure}[H]
\centering
\begin{tikzpicture}[node distance=0cm, auto,
    every node/.style={font=\small},
    block/.append style={minimum width=1.6cm},
    purpleblock/.append style={minimum width=1.4cm}
]

    % Main pipeline nodes
    \node (input)   []                              {Input $x$};
    \node (qkv)     [block, right=0.6cm of input]   {QKV Proj.\\(Conv 1x1)};

    % Q, K, V stacked vertically
    \node (k)       [purpleblock, right=0.8cm of qkv]               {Key $K$};
    \node (q)       [purpleblock, above=0.2cm of k]                 {Query $Q$};
    \node (v)       [purpleblock, below=0.2cm of k]                 {Value $V$};

    % Softmax block
    \node (softmax) [grayblock, right=0.8cm of k, yshift=0.5cm, minimum width=2.2cm, font=\footnotesize]
                    {Softmax\\$\left(\dfrac{Q{\cdot}K^T}{\sqrt{d_k}}\right)$};

    % Attn*V block on the main line
    \node (matmul)  [grayblock, right=3.8cm of k]                   {Attn$\,{\cdot}\,V$};

    % PE node below matmul
    \node (pe)      [block, below=0.6cm of matmul]                  {PE (DWConv 3x3)};

    % Sum nodes on the main line
    \node (sum_pe)  [sum, right=0.8cm of matmul];
    \node (sum_res) [sum, right=0.8cm of sum_pe];
    \node (output)  [right=0.6cm of sum_res]                        {Output};

    % --- connections ---
    % Input to QKV
    \path [line] (input) -- (qkv);

    % QKV fan-out to Q, K, V
    \path [line] (qkv.east) -- ++(0.4,0) |- (q.west);
    \path [line] (qkv.east) -- ++(0.4,0) |- (k.west);
    \path [line] (qkv.east) -- ++(0.4,0) |- (v.west);

    % Q and K to Softmax (entering separate vertical offsets horizontally, no overlaps)
    % Q.east enters Softmax at y = +0.25 (relative to softmax center)
    % K.east enters Softmax at y = -0.25 (relative to softmax center)
    \coordinate (Q_in) at ($(softmax.west)+(0, 0.25)$);
    \coordinate (K_in) at ($(softmax.west)+(0,-0.25)$);
    
    \path [line] (q.east) -- ++(0.4,0) |- (Q_in);
    \path [line] (k.east) -- ++(0.4,0) |- (K_in);

    % Softmax to Attn*V
    \path [line] (softmax.east) -- ++(0.4,0) |- (matmul.west);

    % V to Attn*V (south port) and PE (west port)
    % Single horizontal segment from V east, splitting at vbranch
    \coordinate (vbranch) at ($(v.east)+(0.4,0)$);
    \coordinate (v_up)    at ($(vbranch)+(0,0.35)$);
    \draw [-] (v.east) -- (vbranch);
    
    % Branch to PE: down then horizontally into PE west
    \path [line] (vbranch) |- (pe.west);
    
    % Branch to Attn*V: go up to clear PE, then horizontally right, then straight up into matmul south
    \path [line] (vbranch) -- (v_up) -- (v_up -| matmul.south) -- (matmul.south);

    % Attn*V and PE converge at sum_pe (entering north/south ports vertically)
    \path [line] (matmul.east) -| (sum_pe.north);
    \path [line] (pe.east)     -| (sum_pe.south);

    % sum_pe to sum_res to output
    \path [line] (sum_pe)  -- (sum_res);
    \path [line] (sum_res) -- (output);

    % Global residual connection (centered label on horizontal wire)
    \coordinate (Rup) at ($(input.north)+(0, 1.9)$);
    \path [line] (input.north) -- (Rup) -- node[above, midway, font=\small\itshape] {Global Residual Connection} (Rup -| sum_res.north) -- (sum_res.north);
\end{tikzpicture}
\caption{Sơ đồ khối Attention với QKV projection và Position Encoding.}
\label{fig:attention}
\end{figure}

\begin{align}
    Q, K, V &= \text{split}(\text{QKV\_Conv}(x)) \\
    \text{Attn} &= \text{softmax}\left(\frac{Q \cdot K^T}{\sqrt{d_k}}\right) \\
    \text{Output} &= V \cdot \text{Attn}^T + \text{PE}(V)
\end{align}

\begin{itemize}
    \item \textbf{QKV Conv}: Lớp tích chập Conv $1\times1$ chiếu đồng thời vector đầu vào sang 3 không gian Query ($Q$), Key ($K$), Value ($V$). Lớp này không sử dụng activation.
    \item $d_k = \texttt{head\_dim} \times \texttt{attn\_ratio}$: Kích thước vector Key được nén nhỏ để tiết kiệm bộ nhớ.
    \item \textbf{PE (Position Encoding)}: Sử dụng một khối Depthwise Conv $3\times3$ tác động lên Value ($V$) nhằm bảo toàn thông tin cấu trúc không gian hình học vốn bị mất trong cơ chế tính toán ma trận Attention thông thường.
    \item \textbf{Multi-head}: Kênh được chia thành $h$ heads riêng biệt để xử lý song song các vùng đặc trưng độc lập.
\end{itemize}

\begin{figure}[H]
\centering
\begin{tikzpicture}[node distance=1.0cm and 1.2cm, auto]
    % Nodes
    \node (x) [grayblock] {Input $x$\\$[B, C, H, W]$};
    \node (qkv) [block, right=of x] {QKV Projection\\(Conv 1x1)};
    \node (split) [purpleblock, right=of qkv] {Query $Q$, Key $K$, Value $V$\\$[B, h, N, d]$};
    \node (attn) [purpleblock, right=of split] {Attention Matrix $A$\\$[B, h, N, N]$};
    \node (out_v) [purpleblock, below=of attn, yshift=0.3cm] {Output Attention\\$[B, h, N, d]$};
    \node (reshape) [block, left=of out_v] {Reshape \& Projection};
    \node (out) [grayblock, left=of reshape] {Output tensor\\$[B, C, H, W]$};

    % Connections
    \path [line] (x) -- (qkv);
    \path [line] (qkv) -- (split);
    \path [line] (split) -- node[above, font=\small] {$Q \cdot K^T$} (attn);
    \path [line] (attn) -- node[right, font=\small] {$\text{Attn} \cdot V$} (out_v);
    \path [line] (out_v) -- (reshape);
    \path [line] (reshape) -- (out);
\end{tikzpicture}
\caption{Chi tiết luồng xử lý và sự biến đổi kích thước tensor trong module Attention.}
\label{fig:attention_flow}
\end{figure}

\subsection{Khối PSABlock (Position-Sensitive Attention Block)}
\label{subsec:psablock}

PSABlock tích hợp cơ chế Self-Attention bên trên và một khối Feed-Forward Network (FFN) theo dạng nối tiếp, đi kèm hai kết nối tắt (residual connection):

\begin{figure}[H]
\centering
\begin{tikzpicture}[node distance=0.8cm and 1.2cm, auto]
    % Nodes
    \node (input) [splitnode] {};
    \node (lbl_in) [left=0.15cm of input] {Input $x$};
    \node (attn) [block, right=of input] {Attention};
    \node (sum1) [sum, right=of attn];
    \node (split1) [splitnode, right=of sum1, xshift=-0.3cm] {};
    \node (ffn) [orangeblock, right=of split1] {FFN};
    \node (sum2) [sum, right=of ffn];
    \node (output) [right=0.8cm of sum2] {Output};

    % Connections
    \path [line] (input) -- (attn);
    \path [line] (attn) -- (sum1);
    \path [line] (sum1) -- (split1);
    \path [line] (split1) -- (ffn);
    \path [line] (ffn) -- (sum2);
    \path [line] (sum2) -- (output);
    
    % Skip 1 (Attention residual)
    \path [line] (input) -- ++(0,1.0) -| (sum1.north);
    
    % Skip 2 (FFN residual)
    \path [line] (split1) -- ++(0,1.0) -| (sum2.north);
\end{tikzpicture}
\caption{Sơ đồ PSABlock: Attention + FFN với các kết nối residual.}
\label{fig:psablock}
\end{figure}

\begin{align}
    x_1 &= x + \text{Attention}(x) \quad \text{(Residual Attention)} \\
    x_2 &= x_1 + \text{FFN}(x_1) \quad \text{(Residual FFN)}
\end{align}

FFN (Feed-Forward Network) gồm lớp Conv $1\times1$ tăng số kênh ($c \to 2c$), đi qua hàm kích hoạt SiLU, và lớp Conv $1\times1$ nén kênh ngược về ban đầu ($2c \to c$, act=False):
\begin{equation}
    \text{FFN}(x_1) = \text{cv2}\Big(\text{SiLU}\big(\text{cv1}(x_1)\big)\Big)
\end{equation}

\begin{figure}[H]
\centering
\begin{tikzpicture}[node distance=1.0cm and 1.2cm, auto]
    % Nodes
    \node (input) [] {Input $x_1$ (Kênh $c$)};
    \node (cv1) [block, right=of input] {Conv 1x1 + SiLU\\(cv1: $c \to 2c$)};
    \node (cv2) [block, right=of cv1] {Conv 1x1\\(cv2: $2c \to c$)};
    \node (output) [right=of cv2] {Output (Kênh $c$)};

    % Connections
    \path [line] (input) -- (cv1);
    \path [line] (cv1) -- (cv2);
    \path [line] (cv2) -- (output);
\end{tikzpicture}
\caption{Sơ đồ khối FFN (Feed-Forward Network) trong PSABlock.}
\label{fig:ffn_block}
\end{figure}

\subsection{Khối C2PSA (Channel-Split Position-Sensitive Attention)}
\label{subsec:c2psa}

Khối C2PSA áp dụng cơ chế chia nhánh (Channel-Split) của CSPNet kết hợp với PSABlock nhằm tối ưu hóa chi phí tính toán của cơ chế Self-Attention:

\begin{figure}[H]
\centering
\begin{tikzpicture}[node distance=0.6cm and 0.8cm, auto]
    % Nodes
    \node (input) [] {Input $x$};
    \node (cv1) [block, right=of input] {Conv 1x1\\(cv1)};
    \node (split) [splitnode, right=of cv1, xshift=-0.2cm] {};
    \node (psablock) [orangeblock, right=of split] {PSABlock};
    \node (concat) [grayblock, right=of psablock] {Concat};
    \node (cv2) [block, right=of concat] {Conv 1x1\\(cv2)};
    \node (output) [right=of cv2] {Output};

    % Connections
    \path [line] (input) -- (cv1);
    \path [line] (cv1) -- (split);
    \path [line] (split) -- (psablock);
    \path [line] (psablock) -- (concat);
    \path [line] (concat) -- (cv2);
    \path [line] (cv2) -- (output);
    
    % Upper shortcut (Branch a)
    \coordinate (SplitUp) at ($(split)+(0,1.0)$);
    \path [line] (split) -- (SplitUp) -- node[above, midway, font=\small] {Nhánh $a$ (đi thẳng)} (SplitUp -| concat.north) -- (concat.north);
\end{tikzpicture}
\caption{Sơ đồ C2PSA: CSP splitting + PSABlock.}
\label{fig:c2psa}
\end{figure}

\begin{align}
    a, b &= \text{split}(\text{cv1}(x)) \\
    b' &= \text{PSABlock}(b) \\
    \text{Output} &= \text{cv2}(\text{Concat}[a, b'])
\end{align}

Đầu vào qua \texttt{cv1} được phân đôi: nhánh $a$ đi thẳng bảo toàn đặc trưng tích chập cục bộ, nhánh $b$ đi qua PSABlock học các phụ thuộc không gian xa. Cuối cùng, hai nhánh được concat và hợp nhất qua \texttt{cv2}.

\begin{figure}[H]
\centering
\begin{tikzpicture}[auto]
    % Input grid
    \draw[step=0.2cm, gray!20, very thin] (0,0) grid (2,2);
    \draw[thick, draw=black] (0,0) rectangle (2,2);
    \draw[fill=blue!40] (0.8, 0.8) rectangle (1.2, 1.2);
    \node[below=0.3cm, font=\small\bfseries, align=center] at (1, 0) {Đặc trưng cục bộ};
    
    % Arrow 1
    \draw[thick, ->] (2.4, 1) -- node[above=2pt, font=\small\bfseries] {SPPF} (3.9, 1);
    
    % SPPF output grid (shifted by 4.3cm)
    \begin{scope}[xshift=4.3cm]
        \draw[step=0.2cm, gray!20, very thin] (0,0) grid (2,2);
        \draw[thick, draw=black] (0,0) rectangle (2,2);
        \draw[fill=blue!40] (0.8, 0.8) rectangle (1.2, 1.2);
        % Concentric rings of receptive fields
        \draw[orange, thick, dashed] (1,1) circle (0.4cm);
        \draw[red, thick, dashed] (1,1) circle (0.8cm);
        \node[below=0.3cm, font=\small\bfseries, align=center] at (1, 0) {Đa quy mô\\(Multi-scale)};
    \end{scope}
    
    % Arrow 2
    \draw[thick, ->] (6.7, 1) -- node[above=2pt, font=\small\bfseries] {C2PSA} (8.2, 1);
    
    % C2PSA output grid (shifted by 8.6cm)
    \begin{scope}[xshift=8.6cm]
        \draw[step=0.2cm, gray!20, very thin] (0,0) grid (2,2);
        \draw[thick, draw=black] (0,0) rectangle (2,2);
        \draw[fill=blue!40] (0.8, 0.8) rectangle (1.2, 1.2);
        % Arrows pointing everywhere representing global attention
        \foreach \x in {0.2, 0.6, 1.4, 1.8} {
            \foreach \y in {0.2, 0.6, 1.4, 1.8} {
                \draw[blue!60, -latex, thin] (1,1) -- (\x, \y);
            }
        }
        \node[below=0.3cm, font=\small\bfseries, align=center] at (1, 0) {Ngữ cảnh toàn cục};
    \end{scope}
\end{tikzpicture}
\caption{Mô phỏng sự biến đổi của feature map khi đi qua SPPF và C2PSA.}
\label{fig:sppf_c2psa_transform}
\end{figure}

\begin{notebox}[Tại sao chỉ dùng Attention ở P5?]
Độ phức tạp tính toán của cơ chế Self-Attention là $O(N^2)$ với $N = H \times W$ là số điểm ảnh. Ở tầng P5 ($10 \times 10$), $N=100$ nên chi phí cực kỳ nhẹ. Ngược lại, ở tầng P3 ($40 \times 40$), $N=1600$ khiến độ phức tạp tăng gấp 256 lần, quá nặng cho các mô hình nhỏ như YOLO26-Nano.
\end{notebox}

\subsection{Tổng quan đồ thị các khối đặc biệt}
\label{subsec:blocks_overview}

Để trực quan hóa sự sắp đặt các khối tùy biến, đồ thị dưới đây biểu diễn mối liên kết giữa C3k2, SPPF và C2PSA trong mô hình YOLO26:

\begin{figure}[H]
\centering
\begin{tikzpicture}[node distance=1.0cm and 1.2cm, auto]
    % Nodes
    \node (p5_in) [grayblock] {Đặc trưng P5\\ Backbone};
    \node (c3k2) [block, right=of p5_in] {C3k2 (c3k=True)\\ {[Trích xuất đặc trưng]}};
    \node (sppf) [orangeblock, right=of c3k2] {SPPF\\ {[Mở rộng trường cảm]}};
    \node (c2psa) [purpleblock, right=of sppf] {C2PSA\\ {[Học ngữ cảnh toàn cục]}};
    \node (neck) [grayblock, right=of c2psa] {Neck / Head\\ {[Khối phát hiện]}};

    % Connections
    \path [line] (p5_in) -- (c3k2);
    \path [line] (c3k2) -- (sppf);
    \path [line] (sppf) -- (c2psa);
    \path [line] (c2psa) -- (neck);
\end{tikzpicture}
\caption{Tổng quan cách tổ chức các khối đặc biệt: C3k2, SPPF, và C2PSA trong YOLO26.}
\label{fig:blocks_overview}
\end{figure}

\subsection{Bảng chi tiết các lớp Backbone}

\begin{table}[H]
\centering
\small
\caption{Chi tiết các lớp trong Backbone của YOLO26n (scale nano, ảnh $320\times320$).}
\label{tab:backbone_layers}
\begin{tabular}{clllll}
\toprule
\textbf{Layer} & \textbf{Module} & \textbf{Input} & \textbf{Output (n)} & \textbf{Key Args} & \textbf{Note} \\
\midrule
b0  & Conv   & $3\times320^2$ & $16\times160^2$ & $k\!=\!3, s\!=\!2$ & P1/2 \\
b1  & Conv   & $16\times160^2$ & $32\times80^2$ & $k\!=\!3, s\!=\!2$ & P2/4 \\
b2  & C3k2   & $32\times80^2$ & $64\times80^2$ & c3k=F, $e\!=\!0.25$ & \\
b3  & Conv   & $64\times80^2$ & $64\times40^2$ & $k\!=\!3, s\!=\!2$ & \textbf{P3/8} \\
b4  & C3k2   & $64\times40^2$ & $128\times40^2$ & c3k=F, $e\!=\!0.25$ & \\
b5  & Conv   & $128\times40^2$ & $128\times20^2$ & $k\!=\!3, s\!=\!2$ & \textbf{P4/16} \\
b6  & C3k2   & $128\times20^2$ & $128\times20^2$ & c3k=T & \\
b7  & Conv   & $128\times20^2$ & $256\times10^2$ & $k\!=\!3, s\!=\!2$ & \textbf{P5/32} \\
b8  & C3k2   & $256\times10^2$ & $256\times10^2$ & c3k=T & \\
b9  & SPPF   & $256\times10^2$ & $256\times10^2$ & $k\!=\!5, n\!=\!3$ & \\
b10 & C2PSA  & $256\times10^2$ & $256\times10^2$ & attention & \\
\bottomrule
\end{tabular}
\end{table}

% ===================================================================
\section{Neck}
\label{sec:arch_neck}

Neck của \yolo sử dụng kiến trúc kết hợp FPN (top-down) và PAN (bottom-up) để tổng hợp đặc trưng đa tỷ lệ từ Backbone trước khi đưa vào Detection Head.

\subsection{Top-Down Pathway (FPN)}

Truyền semantic information từ tầng sâu (P5) xuống tầng nông (P3):

\begin{enumerate}
    \item Upsample b10 ($256 \times 10^2$) $\times 2 \to 256 \times 20^2$
    \item Concat với b6 ($128 \times 20^2$) $\to 384 \times 20^2$
    \item C3k2 $\to$ h13: $128 \times 20^2$
    \item Upsample h13 $\times 2 \to 128 \times 40^2$
    \item Concat với b4 ($128 \times 40^2$) $\to 256 \times 40^2$
    \item C3k2 $\to$ h16: $64 \times 40^2$ \quad (P3 output)
\end{enumerate}

\subsection{Bottom-Up Pathway (PAN)}

Truyền detail information từ tầng nông (P3) lên tầng sâu (P5):

\begin{enumerate}
    \item Conv stride-2: h16 ($64 \times 40^2$) $\to 64 \times 20^2$
    \item Concat với h13 ($128 \times 20^2$) $\to 192 \times 20^2$
    \item C3k2 $\to$ h19: $128 \times 20^2$ \quad (P4 output)
    \item Conv stride-2: h19 ($128 \times 20^2$) $\to 128 \times 10^2$
    \item Concat với b10 ($256 \times 10^2$) $\to 384 \times 10^2$
    \item C3k2(attn=True) $\to$ h22: $256 \times 10^2$ \quad (P5 output)
\end{enumerate}

Neck tạo ra 3 đầu ra ở 3 mức phân giải:
\begin{itemize}
    \item \textbf{P3\_out} (h16): $80 \times 80$, stride 8 --- phát hiện vật thể \textbf{nhỏ}.
    \item \textbf{P4\_out} (h19): $40 \times 40$, stride 16 --- phát hiện vật thể \textbf{trung bình}.
    \item \textbf{P5\_out} (h22): $20 \times 20$, stride 32 --- phát hiện vật thể \textbf{lớn}.
\end{itemize}

\begin{table}[H]
\centering
\caption{Số anchor points tại mỗi scale detection (ảnh $320\times320$).}
\label{tab:anchors}
\begin{tabular}{lcccc}
\toprule
\textbf{Scale} & \textbf{Stride} & \textbf{Feature Map} & \textbf{Anchor Points} & \textbf{Phù hợp} \\
\midrule
P3 & 8  & $40\times40$  & 1600 & Vật nhỏ \\
P4 & 16 & $20\times20$  & 400  & Vật vừa \\
P5 & 32 & $10\times10$  & 100  & Vật lớn \\
\midrule
\multicolumn{3}{r}{\textbf{Tổng}} & \textbf{2100} & \\
\bottomrule
\end{tabular}
\end{table}

% ===================================================================
\section{Detection Head: Dual-Head O2M/O2O}
\label{sec:arch_head}

\yolo sử dụng kiến trúc Dual-Head với hai nhánh dự đoán song song:

\subsection{Nhánh cv2 --- Box Prediction}

Dự đoán tọa độ bounding box thông qua 2 lớp Conv $3 \times 3$ và một lớp Conv $1 \times 1$ cuối:

\begin{lstlisting}[caption={Nhánh cv2 (Box) trong Detection Head}, label=lst:cv2]
self.cv2 = nn.ModuleList(
    nn.Sequential(
        Conv(x, c2, 3),                     # Conv 3x3
        Conv(c2, c2, 3),                    # Conv 3x3
        nn.Conv2d(c2, 4 * self.reg_max, 1)  # Conv 1x1 -> 4*reg_max
    ) for x in ch
)
\end{lstlisting}

Đầu ra: $4 \times \text{reg\_max}$ kênh. Với Direct Regression ($\text{reg\_max} = 1$), đầu ra chỉ có 4 kênh tương ứng 4 khoảng cách LTRB (left, top, right, bottom).

\subsection{Nhánh cv3 --- Class Prediction}

Dự đoán xác suất lớp đối tượng sử dụng Depthwise Conv để tối ưu tốc độ:

\begin{lstlisting}[caption={Nhánh cv3 (Class) trong Detection Head}, label=lst:cv3]
self.cv3 = nn.ModuleList(
    nn.Sequential(
        nn.Sequential(DWConv(x, x, 3), Conv(x, c3, 1)),  # DWConv + Conv
        nn.Sequential(DWConv(c3, c3, 3), Conv(c3, c3, 1)),
        nn.Conv2d(c3, self.nc, 1)  # Conv 1x1 -> nc classes
    ) for x in ch
)
\end{lstlisting}

\subsection{Cơ chế Dual-Head}

\begin{table}[H]
    \centering
    \caption{So sánh nhánh O2M và O2O trong Dual-Head}
    \label{tab:dual_head}
    \begin{tabularx}{\textwidth}{lXX}
        \toprule
        \textbf{Đặc tính} & \textbf{O2M (One-to-Many)} & \textbf{O2O (One-to-One)} \\
        \midrule
        Khi nào dùng & Chỉ khi \textbf{huấn luyện} & Khi \textbf{suy luận} (và cả huấn luyện) \\
        TAL topk     & $\text{topk} = 10$ & $\text{topk}_2 = 1$ \\
        Số dự đoán / vật thể & Nhiều anchor $\to$ 1 vật thể & Đúng 1 anchor $\to$ 1 vật thể \\
        Mục đích     & Tạo gradient dày đặc, đẩy nhanh tốc độ học & NMS-free, dự đoán sạch \\
        Cần NMS?     & Có (khi suy luận đơn lẻ) & \textbf{Không} \\
        \bottomrule
    \end{tabularx}
\end{table}

Khi deploy thực tế, nhánh O2M bị loại bỏ hoàn toàn thông qua hàm \texttt{strip\_o2m()}, chỉ giữ lại nhánh O2O phục vụ suy luận trực tiếp.

\begin{table}[H]
\centering
\caption{So sánh model trước và sau khi strip O2M.}
\label{tab:strip_o2m}
\begin{tabular}{lccc}
\toprule
\textbf{Metric} & \textbf{Training (last.pt)} & \textbf{Inference (lite.pt)} & \textbf{Giảm} \\
\midrule
File size & 38.2 MB & 5.6 MB & $\sim$85\% \\
Precision & FP32 & FP16 & 50\% \\
O2M Head & Có & Stripped & -- \\
Optimizer state & Có & Stripped & -- \\
\bottomrule
\end{tabular}
\end{table}

% ===================================================================
\section{Box Decoding: Direct Regression}
\label{sec:arch_box_decoding}

\yolo sử dụng phương pháp hồi quy trực tiếp (Direct Regression) thay cho DFL~\cite{dfl}, giúp giảm đáng kể khối lượng tính toán:

\begin{enumerate}[label=\textbf{Bước \arabic*:}]
    \item \textbf{Feature Extraction}: Trích xuất đặc trưng qua Backbone và Neck.
    \item \textbf{Tạo Anchor Grids}: Xác định tọa độ tâm $(c_x, c_y)$ cho mỗi ô lưới.
    \item \textbf{Dự đoán LTRB}: Nhánh cv2 xuất ra 4 giá trị khoảng cách: $d_\text{left}$, $d_\text{top}$, $d_\text{right}$, $d_\text{bottom}$.
    \item \textbf{Khôi phục tọa độ}:
    \begin{equation}
        x_\text{min} = c_x - d_\text{left}, \quad y_\text{min} = c_y - d_\text{top}
    \end{equation}
    \begin{equation}
        x_\text{max} = c_x + d_\text{right}, \quad y_\text{max} = c_y + d_\text{bottom}
    \end{equation}
    \item \textbf{Rescale}: Nhân tọa độ với stride tương ứng để đưa về kích thước ảnh gốc.
\end{enumerate}

\subsection{Cơ chế DFL (Distribution Focal Loss)}
\label{subsec:dfl_module}

DFL~\cite{dfl} chuyển bài toán hồi quy khoảng cách liên tục thành bài toán phân loại trên một số lượng $\texttt{reg\_max}$ các bins rời rạc:

\begin{equation}
    \hat{d} = \sum_{i=0}^{\texttt{reg\_max}-1} i \cdot p_i, \quad p_i = \text{softmax}(\text{logits})_i
\end{equation}

Trọng số của lớp Conv $1\times1$ cuối được cố định là chuỗi tăng dần $[0, 1, 2, \ldots, \texttt{reg\_max}-1]$ và không tham gia vào quá trình lan truyền ngược để cập nhật tham số.

\begin{notebox}[YOLO26 mặc định: reg\_max=1 --- không dùng DFL]
Khi đặt cấu hình mặc định \texttt{reg\_max=1}, DFL sẽ tự động được thay thế bằng một lớp đồng nhất \texttt{nn.Identity()}. Lúc này mô hình thực hiện hồi quy trực tiếp 4 giá trị khoảng cách $(l, t, r, b)$, đầu ra của box head chỉ gồm 4 kênh. Với cấu hình \texttt{reg\_max=16}, đầu ra sẽ có 64 kênh, cho phép dự đoán phân phối xác suất biên tốt hơn nhưng tăng chi phí tính toán.
\end{notebox}

\subsection{So sánh Direct Regression và DFL}

\begin{table}[H]
    \centering
    \caption{So sánh hồi quy trực tiếp (Direct Regression) và DFL}
    \label{tab:direct_vs_dfl}
    \begin{tabularx}{\textwidth}{lXX}
        \toprule
        \textbf{Đặc tính} & \textbf{Direct Regression} & \textbf{DFL} \\
        \midrule
        Cách tính & Xuất trực tiếp 4 giá trị $l, r, t, b$ & Softmax 16 giá trị/hướng, tổng trọng số \\
        Kênh đầu ra & $4 \times 1 = 4$ & $4 \times 16 = 64$ \\
        Ưu điểm & Tốc độ nhanh, đơn giản & Chính xác hơn ở biên mờ \\
        Nhược điểm & Kém chính xác ở biên mờ & Softmax tốn CPU, gây nghẽn \\
        \bottomrule
    \end{tabularx}
\end{table}

\yolo bù đắp sự suy giảm độ chính xác của Direct Regression bằng cách đặt $\text{topk}=10$ ở nhánh O2M, tạo gradient dày đặc ép Backbone và Neck học đặc trưng sắc nét hơn.

% ===================================================================
\section{Phân tỷ lệ mô hình (Model Scaling)}
\label{sec:arch_scaling}

\yolo hỗ trợ nhiều biến thể kích thước thông qua 3 hệ số:

\begin{table}[H]
\centering
\caption{Cấu hình Scale của YOLO26.}
\label{tab:scales}
\begin{tabular}{lcccl}
\toprule
\textbf{Scale} & \textbf{Depth} & \textbf{Width} & \textbf{Max Ch} & \textbf{Ghi chú} \\
\midrule
n (nano)   & 0.50 & 0.25 & 1024 & Nhỏ nhất, nhanh nhất \\
s (small)  & 0.50 & 0.50 & 1024 & Cân bằng tốc độ/chính xác \\
m (medium) & 0.50 & 1.00 & 512  & Trung bình \\
l (large)  & 1.00 & 1.00 & 512  & Lớn \\
x (xlarge) & 1.00 & 1.50 & 512  & Chính xác nhất \\
\bottomrule
\end{tabular}
\end{table}

\textbf{Công thức tính kênh và lặp:}
\begin{align}
    \texttt{ch}(c) &= \texttt{make\_divisible}\left(\min(c, \texttt{max\_ch}) \times \texttt{width},\; 8\right) \label{eq:ch} \\
    \texttt{rep}(n) &= \max\left(\text{round}(n \times \texttt{depth}),\; 1\right) \label{eq:rep}
\end{align}

Ví dụ scale nano: $\texttt{ch}(256) = \texttt{make\_divisible}(\min(256, 1024) \times 0.25, 8) = 64$.

Trong đồ án, nhóm sử dụng biến thể \textbf{nano} (n) với tổng dung lượng weights khoảng 5.3MB, phù hợp cho triển khai trên \rpi.
